% DRAFT — apply during W4 restructure (replaces paper/sections/0_abstract.tex)
% Conditional on DiLiGenT-MV producing a positive captured-target signal
% (per the codex W1 D7 decision gate). If captured is negative, fall back
% to the current 7.5 abstract.

UV atlases for procedural, generated, and noisy assets are usually
judged by geometry-facing proxies, but production PBR pipelines fail
through a different interface: residual hot spots, mip leakage, and
cross-channel seams that only become visible after baking and
relighting. We introduce a residual-driven selective UV atlas repair
method that treats the held-out PBR relighting residual as the
deployment signal. The pipeline bakes albedo, normal, roughness, and
metallic channels through a differentiable PBR renderer, attributes
residuals back to faces and charts, repairs only high-residual local
charts, reallocates a fixed texel budget with mip-aware demand, and
deploys the candidate atlas only if a matched-control C5 gate accepts
it on a held-out gate split. We validate the method against three
target-signal regimes: a controlled synthetic oracle (spot, bunny,
Objaverse), a self-described oracle on procedural noisy and
SDS-pipeline init shapes, and \emph{captured imagery} from
DiLiGenT-MV (5 objects, 20 views, 96 calibrated lights/view), with
disjoint proposal/gate/test splits and multi-seed reruns. Across
captured DiLiGenT-MV cases the gate triggers on the residual-heavy
PartUV atlases while rolling back already-strong xatlas baselines;
across the synthetic/self-described regimes the same selectivity
holds (3/20 main accepts, 2/5 SDS-init real meshes). A strict
matched-control study shows the gain survives atlas-size, padding,
and chart-count locks but vanishes under distortion-tail or
utilization locks. The method is therefore not a universal UV
optimizer; it is a selective, rollback-gated way to trade localized
UV proxy metrics for substantially better downstream PBR fidelity,
validated against \emph{captured} as well as oracle target signals.
